\section{结论}\label{sec:conclusion}
本文提出 EGP-NPA，将等变残差策略、局部动态份额分配与完成包络结合用于耦合安全约束下的多机器人重排。七机器人证书揭示了误差单调控制在安全可达任务中的局限，两节点子问题说明了动态份额转移产生协同退让的机制，完成包络则保证残差一阶消隐并恢复目标零速平衡点。

一次完整预算训练的冻结策略在 $96$ 个开发任务及 $48$ 个留出任务上均安全到达。进一步的八次匹配预算训练和 $384$ 次评测表明，投影正则对到达数的影响随价值网络配置改变：团队价值网络下减少到达，局部价值网络下增加到达。静态分配结果作为固定权重消融报告，认证和计算预算终止与零速停滞分别保留。另以官方 DGPPO 更新开展了两次预算受控训练和新任务评测，结果及适配差异见表\ref{tab:external}；内部消融不再作为外部算法的替代。

目前学习变体每格只有两个训练重复；主单积分器评测的谱特征来自精确谱计算，二阶物理实验则另行接入在线估计器。图策略每层沿一跳边传递消息，多层推理具有更大的信息作用范围。物理实验中既有距离越界，也有保守补偿下的未到达，表明执行器和通信条件需纳入更完整的闭环分析。后续独立方法比较应采用对应原算法的实现、共同任务适配、可比总交互与调参预算、更多训练重复及新的留出任务。
